\chapter{物理常数表}

本附录列出高中至大学低年级物理中最常用的常数、单位与天体数据。数值取常用近似值，足以满足教材中的数量级估算与典型题计算。

\section{基本物理常数}

\begin{longtable}{p{0.22\textwidth}p{0.18\textwidth}p{0.22\textwidth}p{0.26\textwidth}}
\toprule
常数名称 & 符号 & 数值 & 常用单位 \\
\midrule
\endfirsthead

\toprule
常数名称 & 符号 & 数值 & 常用单位 \\
\midrule
\endhead

\bottomrule
\endfoot

真空光速 & $c$ & $2.99792458\times 10^8$ & $\text{m}\,\text{s}^{-1}$ \\
万有引力常量 & $G$ & $6.67430\times 10^{-11}$ & $\text{N}\,\text{m}^2\,\text{kg}^{-2}$ \\
普朗克常量 & $h$ & $6.62607015\times 10^{-34}$ & $\text{J}\,\text{s}$ \\
元电荷 & $e$ & $1.602176634\times 10^{-19}$ & $\text{C}$ \\
电子质量 & $m_e$ & $9.1093837\times 10^{-31}$ & $\text{kg}$ \\
质子质量 & $m_p$ & $1.6726219\times 10^{-27}$ & $\text{kg}$ \\
玻尔兹曼常量 & $k_B$ & $1.380649\times 10^{-23}$ & $\text{J}\,\text{K}^{-1}$ \\
真空介电常量 & $\varepsilon_0$ & $8.8541878128\times 10^{-12}$ & $\text{F}\,\text{m}^{-1}$ \\
真空磁导率 & $\mu_0$ & $4\pi\times 10^{-7}$ & $\text{N}\,\text{A}^{-2}$ \\
阿伏伽德罗常量 & $N_A$ & $6.02214076\times 10^{23}$ & $\text{mol}^{-1}$ \\
\end{longtable}

\section{SI 基本单位}

\begin{longtable}{p{0.24\textwidth}p{0.14\textwidth}p{0.48\textwidth}}
\toprule
物理量 & 单位符号 & 单位名称 \\
\midrule
\endfirsthead

\toprule
物理量 & 单位符号 & 单位名称 \\
\midrule
\endhead

\bottomrule
\endfoot

长度 & m & 米 \\
质量 & kg & 千克 \\
时间 & s & 秒 \\
电流 & A & 安培 \\
热力学温度 & K & 开尔文 \\
物质的量 & mol & 摩尔 \\
发光强度 & cd & 坎德拉 \\
\end{longtable}

\section{常用换算}

\begin{longtable}{p{0.42\textwidth}p{0.42\textwidth}}
\toprule
换算关系 & 说明 \\
\midrule
\endfirsthead

\toprule
换算关系 & 说明 \\
\midrule
\endhead

\bottomrule
\endfoot

$1\,\text{eV}=1.602176634\times 10^{-19}\,\text{J}$ & 微观粒子能量常用单位 \\
$1\,\text{u}=1.66053906660\times 10^{-27}\,\text{kg}$ & 原子质量单位 \\
$1\,\text{atm}=1.01325\times 10^5\,\text{Pa}$ & 标准大气压 \\
$1\,\text{L}=10^{-3}\,\text{m}^3$ & 体积换算 \\
$1\,\text{cm}^3=1\,\text{mL}$ & 小体积常用换算 \\
$1\,\text{kW}\cdot\text{h}=3.6\times 10^6\,\text{J}$ & 电能计费单位 \\
$0\,^{\circ}\text{C}=273.15\,\text{K}$ & 摄氏温标与绝对温标 \\
$1\,\text{ly}\approx 9.46\times 10^{15}\,\text{m}$ & 光年 \\
$1\,\text{AU}\approx 1.496\times 10^{11}\,\text{m}$ & 天文单位 \\
\end{longtable}

\section{太阳系数据}

\begin{longtable}{p{0.18\textwidth}p{0.26\textwidth}p{0.24\textwidth}p{0.20\textwidth}}
\toprule
天体 & 质量 & 平均半径 & 备注 \\
\midrule
\endfirsthead

\toprule
天体 & 质量 & 平均半径 & 备注 \\
\midrule
\endhead

\bottomrule
\endfoot

太阳 & $1.989\times 10^{30}\,\text{kg}$ & $6.96\times 10^8\,\text{m}$ & 太阳系质量主体 \\
地球 & $5.972\times 10^{24}\,\text{kg}$ & $6.371\times 10^6\,\text{m}$ & 表面重力加速度约 $9.8\,\text{m}\,\text{s}^{-2}$ \\
月球 & $7.35\times 10^{22}\,\text{kg}$ & $1.737\times 10^6\,\text{m}$ & 地月平均距离约 $3.84\times 10^8\,\text{m}$ \\
木星 & $1.898\times 10^{27}\,\text{kg}$ & $6.99\times 10^7\,\text{m}$ & 太阳系最大行星 \\
\end{longtable}

\section{使用说明}

\begin{itemize}
  \item 在考试与课堂练习中，若题目未要求高精度，通常可采用两到三位有效数字。
  \item 估算题中常把 $c\approx 3.0\times 10^8\,\text{m/s}$、$g\approx 9.8\,\text{m/s}^2$、$G\approx 6.67\times 10^{-11}$ 作为常用近似。
  \item 天体运动问题中，质量、半径和轨道半径往往配合万有引力公式与圆周运动公式使用。
\end{itemize}
